\begin{chapintro}
本文から取り出すデータの中心は 2×2 表の 4 つの数で、感度と特異度しか書かれていなくても、人数が分かれば表を復元できます。取り出した研究ごとに、QUADAS-2 の 4 つの領域でバイアスのリスクと適用可能性を判定します。検査の研究に特有のバイアスのうち初めに押さえる 4 つと、検査技師が論文の方法の節で確かめる点を整理し、最後に CSV を R に読み込みます。実証研究の紹介と、2026 年に公表された改訂版の QUADAS-3 の変更点は、発展の印をつけて載せました。
\end{chapintro}

\section{抽出フォームの項目}
\label{sec:12-1}

データ抽出は、\kkref{sec:04-1}と同じく 2 人が独立に行い、結果を突き合わせます。抽出の前に、何をどの形で記録するかを抽出フォームとして決め、2〜3 本の論文で試して直しておきます。検査のレビューで記録する項目は、おおむね次のとおりです。
\begin{description}
\item[研究の特徴] 著者、発表年、国、施設、登録期間、研究デザイン（前向きか後ろ向きか、連続登録か）を記録します。
\item[対象] 場面（救急外来か ICU か）、組み入れ基準と除外基準、年齢、重症度の指標、抗菌薬がすでに投与されていたかを記録します。
\item[対象検査] 測定試薬と機器（測定原理と製造元）、検体の種類（血清か血漿か）、採血の時点（受診時か、発熱から何時間後か）、採血から測定までの時間と保存条件、カットオフとその決め方（事前に決めたか、データを見て決めたか）を記録します。
\item[参照基準] 血液培養のセット数、汚染菌の判断の仕方、培養の結果を判定した人が PCT の値を知っていたかを記録します。
\item[結果] カットオフごとの TP、FP、FN、TN と、判定不能や検体不良で解析から外れた人数を記録します。
\end{description}

検査技師が抽出に加わるとき、対象検査の項目はとくに丁寧に記録してください。\kkref{sec:09-6}で説明したように、測定試薬や採血の時点の違いは研究間のばらつきの原因になりえます。論文の方法の節で測定法が「PCT を測定した」としか書かれていなければ、「記載なし」と記録し、空欄にしません。空欄は「確認していない」のか「書かれていない」のかが区別できないからです。

\section{本文から 2×2 表を作る}
\label{sec:12-2}

論文に 2×2 表がそのまま載っていることは、むしろ少数です。2×2 表の意味は\kkref{sec:09-1}、介入の試験での作り方は\kkref{sec:04-2}で説明しました。多くは感度と特異度、あるいは陽性的中割合と陰性的中割合が、人数と一緒に書かれています。架空の Lopez 2011 の結果の段落が、次のように書かれていたとします。

\begin{KKfbox9ao}[演習の素材（架空）]
登録した 482 人のうち、血液培養が陽性だったのは 62 人（12.9\%）だった。PCT の中央値は、培養陽性群で 1.2 ng/mL、陰性群で 0.21 ng/mL だった。カットオフを 0.5 ng/mL としたときの感度は 56.5\%、特異度は 77.9\% だった。
\end{KKfbox9ao}

この段落から 2×2 表を復元します。病気のある人は 62 人なので、真陽性は 62 × 0.565 で 35.03、四捨五入して 35 人です。偽陰性は残りの 62 − 35 で 27 人になります。病気のない人は 482 − 62 で 420 人なので、真陰性は 420 × 0.779 で 327.18、四捨五入して 327 人、偽陽性は 420 − 327 で 93 人です。復元した表から感度を計算し直すと 35 ÷ 62 で 0.5645、特異度は 327 ÷ 420 で 0.7786 となり、論文の 56.5\% と 77.9\% に一致します。

計算し直して元の値に戻ることを、毎回確かめてください。論文の割合は丸めて書かれているので、人数に掛けた値が 2 つの整数の中間近くに来て、四捨五入の向きを決めかねることがあります。そのときは、その前後の整数で感度を計算し直し、論文の値に一致するほうを選びます。どちらでも一致する場合や、どちらでも一致しない場合は、抽出した 2 人で判断を記録し、必要なら著者に問い合わせます。陽性的中割合と陰性的中割合しか書かれていなくても、陽性になった人数と陰性になった人数が分かれば、同じように表を復元できます。

\section{抽出で迷う場面}
\label{sec:12-3}

2×2 表を作る段階で、判断に迷う場面がいくつかあります。どう扱うかは、データを見る前にプロトコルで決めておきます。研究の結果を見てから扱いを選ぶと、選んだ人の期待に沿った方向へ結果が偏るからです。
\begin{description}
\item[複数のカットオフ] 1 本の論文が 0.25、0.5、1.0 ng/mL のそれぞれで成績を報告していることがあります。同じ患者から 3 つの 2×2 表を取ると、同じ人を 3 回数えることになるので、1 研究から 1 つを選びます。例題では「0.5 ng/mL があればそれを、なければ著者が主とした値を取る」という規則にしました。
\item[判定不能の結果] 検体の溶血や量の不足で測定できなかった人は、2×2 表のどこにも入りません。何人が外れたかを記録し、多ければ\kkref{sec:12-4}で説明するフローとタイミングの領域で考慮します。
\item[1 人の複数のエピソード] 同じ患者が入院中に何度も発熱し、そのたびに PCT と血液培養を測っていることがあります。エピソードを単位に数えると、同じ人の結果が互いに似ているため、独立した観測として扱えません。患者を単位にした表があればそれを取り、なければその旨を記録します。
\item[全員が参照基準を受けていない] 主治医が必要と判断した人だけに血液培養を採った研究では、PCT が低くて培養を採らなかった人が表から抜けます。この偏りを部分的検証バイアスと呼びます。
\end{description}

\section{QUADAS-2 の 4 つの領域}
\label{sec:12-4}

組み入れた研究ごとに、バイアスのリスクを評価します。現在のレビューで最も広く使われているのは QUADAS-2 です（Whiting 2011）。RCT で使う RoB 2 と同じく、領域ごとに判定する道具です\kksee{sec:04-3}。QUADAS-2 は研究を 4 つの領域に分け、領域ごとに「バイアスのリスク」を判定し、最初の 3 つの領域では「適用可能性の懸念」も判定します。バイアスのリスクは、研究の結果が真の精度から系統的にずれている恐れを表し、適用可能性の懸念は、研究の対象や検査の方法がレビューの問いとずれている恐れを表します。
\begin{description}
\item[患者選択] 対象者を連続して、またはランダムに登録したか、症例対照デザインを避けたか、不適切な除外をしていないかを問います。
\item[対象検査] 対象検査の結果を、参照基準の結果を知らずに判定したか、カットオフを事前に決めていたかを問います。
\item[参照基準] 参照基準が標的疾患を正しく判定できる方法か、参照基準の結果を、対象検査の結果を知らずに判定したかを問います。
\item[フローとタイミング] 対象検査と参照基準の間隔は適切か、全員が参照基準を受けたか、全員が同じ参照基準を受けたか、全員が解析に含まれたかを問います。
\end{description}

それぞれの問いはシグナリング質問と呼ばれ、「はい」「いいえ」「不明」で答えます。すべて「はい」なら、その領域のバイアスのリスクは低いと判定します。1 つでも「いいえ」があれば、バイアスが生じる可能性を検討し、その影響が大きいと判断すれば高いと判定します。情報が足りなければ「不明」とします。判定は機械的な点数ではなく、レビュー著者の判断で、判断の理由を一文ずつ記録しておきます。

\section{検査の研究に特有のバイアス}
\label{sec:12-5}

診断精度研究のバイアスは、介入の RCT のバイアスとは別の名前で呼ばれます。QUADAS-2 のシグナリング質問は、これらのバイアスを防げているかを問う形になっています。表 12.1 に、バイアスの名前、何が起きるか、精度が動く向き、対応する QUADAS-2 の領域を並べます。初めて読むときは、上段の 4 つ（スペクトラムバイアス、部分的検証バイアス、診断レビューバイアス、データを見てから決めたカットオフ）を押さえれば十分です。

\begin{kkcenter}
\kkfit{\footnotesize\begin{tabular}{p{24mm}p{43mm}p{29mm}p{16mm}}
\toprule
バイアス & 何が起きるか & 精度の動き & 領域 \\
\midrule
\rowcolor{kkpale}\multicolumn{4}{l}{初めに押さえる 4 つ} \\
スペクトラムバイアス & 重症例と健康者のように、実際の使い方とかけ離れた集団で測る（症例対照デザイン） & 感度も特異度も高く出る & 患者選択 \\
部分的検証バイアス & 一部の人だけが参照基準を受ける。検査陽性の人ほど受けやすい & 感度が高く、特異度が低く出やすい & フローとタイミング \\
診断レビューバイアス & 参照基準の判定者が、対象検査の結果を知っている & 高く出る方向に動くことが多い & 参照基準 \\
データを見てから決めたカットオフ & 最適なカットオフを同じデータで選ぶ & 高く出ると考えられるが、実証は乏しい & 対象検査 \\
\midrule
\rowcolor{kkpale}\multicolumn{4}{l}{そのほか} \\
不適切な除外 & 判定の難しい患者や、境界域の値の患者を解析から除く & 高く出ると考えられる & 患者選択 \\
差異的検証バイアス & 検査の結果によって、別の参照基準を使う & 高く出る方向に動くことが多い & フローとタイミング \\
取り込みバイアス & 対象検査の結果が参照基準の一部になっている & 高く出ると考えられるが、実証は乏しい & 参照基準 \\
検査レビューバイアス & 対象検査の判定者が、参照基準の結果を知っている & 向きは場合による。目視など判定に主観が入る検査で問題になる & 対象検査 \\
対象検査と参照基準の時間差 & 両者の間に病状が変わる（治療の開始、自然経過） & 向きは場合による & フローとタイミング \\
判定不能の結果を除く & 測定できなかった検体や不良検体を解析から除く & 向きは場合による & フローとタイミング \\
\bottomrule
\end{tabular}}
\kkcap{表 12.1　検査の研究に特有のバイアス}
\end{kkcenter}

症例対照デザインの研究は、精度を高く見積もりやすいことが実証研究で示されています。一部の人だけに参照基準を行った研究も、精度が偏りやすくなります。

\begin{hattenbox}
精度が動く向きのうち、実証研究で確かめられているものは限られます。Lijmer らは、診断精度研究のメタ解析で扱われた 184 の研究（218 の検査の評価）を調べ、研究の設計の欠点と診断オッズ比（\kkref{sec:09-4}）の関係を回帰モデルで推定しました（Lijmer 1999）。結果は相対診断オッズ比で示されています。これは、ある欠点をもつ研究の診断オッズ比を、その欠点のない研究の診断オッズ比で割った値で、1 を超えると精度を高く見積もっていることを表します。病気のある集団と別の対照群を比べた研究（症例対照デザイン）は、臨床の集団で調べた研究より診断精度を高く見積もり、相対診断オッズ比は 3.0（95\% 信頼区間 2.0〜4.5）でした。検査の陽性と陰性で別々の参照基準を使った研究も高く見積もり、相対診断オッズ比は 2.2（1.5〜3.3）でした。参照基準を、対象検査の結果を知ったうえで判定した研究も、相対診断オッズ比は 1.3（1.0〜1.9）と高めでした。症例対照デザインで感度も特異度も高く出るのは、病気のある側には見分けやすい重症例が、ない側には明らかに健康な人が集まるためです。

Whiting らの系統的レビューは、研究の結果に影響する要因を、証拠の厚さで分けて報告しています（Whiting 2004）。最も証拠が多かったのは、対象者の特徴、有病割合と重症度、部分的検証バイアス、臨床情報によるレビューバイアス、測定者と測定機器によるばらつきです。対象者の選び方の偏り、参照基準がない、または不適切なこと、差異的検証バイアス、レビューバイアスは、証拠が限られていました。取り込みバイアス、治療パラドックス、閾値の恣意的な選択、脱落については、証拠が見つかりませんでした。表で「実証は乏しい」と書いた取り込みバイアスとデータを見てから決めたカットオフは、この最後の群にあたります。理屈の上では精度が高く出ますが、どれだけ動くかは測られていません。

Rutjes らは、31 件のメタ解析に含まれる 487 研究を使って、15 の研究の特徴と診断精度の関係を調べました（Rutjes 2006）。連続登録でない研究は相対診断オッズ比 1.5（1.0〜2.1）、後ろ向きにデータを集めた研究は 1.6（1.1〜2.2）と、高い精度を報告していました。重症例と健康な対照を組み合わせた研究の推定値は 4.9 ですが、区間は 0.6〜37.3 と広く、確かなことは言えません。この研究の結論は、設計の欠点が精度の推定値に影響しうるが、その大きさは状況によって違う、というものです。表の精度の動きに「場合による」や「実証は乏しい」と書いた項目があるのは、このためです。
\end{hattenbox}

\section{検査技師の目で読む対象検査}
\label{sec:12-6}

対象検査の領域と、その適用可能性の懸念は、検査技師が最も詳しく判断できる部分です。QUADAS-2 のシグナリング質問は、検体の取り扱いや測定条件を直接には問いません。そこで、レビューの問いに合わせて、次の点を論文の方法の節で確かめ、対象検査の領域か適用可能性の懸念に書き込みます。
\begin{description}
\item[検体の取り扱い] 採血から遠心分離までの時間、測定までの時間、保存温度（冷蔵か凍結か）、凍結保存の期間、溶血や乳びの検体の扱いを、方法の節の「検体採取」や「測定」の項で探します。保存した検体を後でまとめて測定した研究は、実際の診療で採血してすぐ測る使い方と条件が違うので、適用可能性の懸念に記録します。何も書かれていなければ「記載なし」とし、判定の理由には「検体の取り扱いの記載がない」と書きます。
\item[測定試薬と機器] 測定原理、製造元、機器の名前、試薬のロットや世代、校正の方法を確かめます。同じ PCT でも試薬が違えば同じ検体で値が少しずれ、カットオフ 0.5 ng/mL の実質的な位置が研究ごとに変わります。レビューの問いが日常検査で使われている自動分析装置を想定していれば、研究用の手法や普及していない試薬を使った研究は、適用可能性の懸念を高とするか、試薬の違いを研究の特徴として記録して\kkref{sec:14-1}の共変量の解析で調べます。
\item[測定者の盲検化] 測定者が血液培養の結果や臨床の経過を知っていたかを、方法の節か、解析に関する記述で確かめます。自動分析装置が出す数値は測定者の判断が入る余地が小さいので、盲検化が書かれていなくても、バイアスのリスクは小さいと考えられます。目視で判定する迅速キットや、波形を人が解釈する検査では、参照基準の結果を知ったうえでの判定が結果を動かしえます。検査の種類によって重みが違うので、判定の理由に検査の種類を書きます。
\item[カットオフの決め方] 方法の節に、カットオフを研究開始前に決めたと明記されているかを確かめます。製造元の添付文書の値や先行研究の値を使っていれば、事前に決めたと判定できます。結果の節で ROC 曲線を描いて最適な点を選んでいれば、データを見て選んだことになり、対象検査の領域は高と判定します。
\end{description}

測定に関わる情報は、本文の方法の節に書かれていないことが多く、補足資料や、研究を計画したときのプロトコル論文にある場合があります。それでも見つからなければ、著者に問い合わせるか、不明と判定します。「PCT を測定した」とだけ書かれた論文を、測定が適切だったと推測して低と判定することは避けます。

\section{QUADAS-2 の演習}
\label{sec:12-7}

架空の 10 研究について、架空の本文の情報から判定した例を図 12.1 に示します。\kkref{sec:04-3}の図 4.1 と同じ形式で、列が領域、行が研究、円の記号が判定です。

\begin{kkcenter}
\begin{minipage}{\linewidth}\centering
\begin{tikzpicture}[x=15mm, y=6.2mm]
\foreach \lab [count=\j] in {患者\\選択, 対象\\検査, 参照\\基準, フローと\\タイミング}
  \node[font=\scriptsize\bfseries, align=center] at (\j+1.0, 0.3) {\lab};
\foreach \name/\a/\b/\c/\d [count=\i] in {
  Kato 2010/l/l/l/l, Lopez 2011/l/l/l/l, Mori 2013/l/l/m/l, Nguyen 2014/l/l/l/m,
  Ono 2015/l/m/l/l, Park 2016/m/l/l/l, Qureshi 2018/l/l/l/h, Rossi 2019/l/h/l/l,
  Sato 2021/l/l/l/l, Tanaka 2023/l/l/l/l} {
  \node[anchor=west, font=\scriptsize] at (-0.9, -\i) {\name};
  \foreach \v [count=\j] in {\a,\b,\c,\d} {
    \if\v l \fill[kksecondary] (\j+1.0,-\i) circle (2.0mm);
      \node[font=\scriptsize] at (\j+1.0,-\i) {+};\fi
    \if\v m \fill[kkpale] (\j+1.0,-\i) circle (2.0mm);
      \draw[kkaccent] (\j+1.0,-\i) circle (2.0mm);
      \node[font=\scriptsize] at (\j+1.0,-\i) {?};\fi
    \if\v h \fill[kkwarm] (\j+1.0,-\i) circle (2.0mm);
      \node[font=\scriptsize, text=white] at (\j+1.0,-\i) {−};\fi
  }
}
\end{tikzpicture}
\kkcap{図 12.1　架空の 10 研究のバイアスのリスクの判定（QUADAS-2）。＋ 低、？ 不明、− 高}
\end{minipage}
\end{kkcenter}

Kato 2010、Lopez 2011、Sato 2021、Tanaka 2023 は、4 つの領域のどれにも懸念がなく、すべて低と判定します。例えば Tanaka 2023 は、救急外来の発熱患者を連続して登録し、カットオフ 0.25 ng/mL を研究開始前の計画書に書いていました。血液培養は全員に採られ、微生物検査の担当者は PCT の値を知らずに判定し、PCT の測定者は培養の結果を知らずに測定していました。

不明の判定は、判定に必要な情報が書かれていない場合につけます。Park 2016 は、患者の登録方法について「研究に同意した患者」としか書かれておらず、連続登録かどうかが分かりません。夜間や休日に受診した重症の患者が抜けていれば、スペクトラムが偏る可能性があります。情報不足なので、患者選択の領域は不明とします。Mori 2013 は、血液培養の陽性と陰性を判定した人が PCT の値を知っていたかどうかの記載がありません。知っていれば診断レビューバイアスが生じうるので、参照基準の領域は不明とします。Nguyen 2014 は、PCT の採血と血液培養の採取の間隔が書かれていません。間隔が長ければ、その間に抗菌薬が始まるなどして菌血症の状態が変わりうるので、フローとタイミングの領域は不明とします。Ono 2015 は、カットオフ 1.0 ng/mL を事前に決めていたのか、結果を見てから選んだのかの記載がありません。対象検査の領域は不明とします。

高の判定は、バイアスが生じる設計であることが分かっている場合につけます。Rossi 2019 は、集めたデータで ROC 曲線を描き、感度と特異度の和が最大になる点としてカットオフ 0.5 ng/mL を選んでいました。データを見てカットオフを選ぶと、その研究の偶然のばらつきに合う位置が選ばれるので、同じカットオフを別の患者に使ったときより精度が高めに出ます。対象検査の領域を高と判定します。

Qureshi 2018 は、ICU の発熱患者のうち、主治医が血液培養を依頼した患者だけを解析していました。PCT が低く、主治医が菌血症を疑わなかった患者は血液培養を採られず、表から抜けています。\kkref{sec:12-3}で説明した部分的検証バイアスにあたり、フローとタイミングの領域を高と判定します。

\section{判定の使い方}
\label{sec:12-8}

QUADAS-2 の判定を点数にして合計し、「質のスコア」として研究を並べることはしません。4 つの領域の偏りは、精度を高める方向にも低める方向にも働き、重みも違うので、足し合わせた数には意味がないからです。判定は領域ごとの表と図（図 12.1 のような信号機の図）で示します。

判定は、研究を解析から除くための関門ではありません。使い道は 3 つあります。1 つ目は、バイアスのリスクが高い研究を除いても結果が変わらないかを確かめる感度分析で、\kkref{sec:14-5}で行います。2 つ目は、判定の違いを共変量やサブグループとして調べる解析です。例題は研究数が少ないので行いません。3 つ目は、結果の確実性を GRADE（\kkref{sec:07-8}）で評価するときの材料で、\kkref{sec:15-2}で使います。

例題では、高と判定された Rossi 2019 と Qureshi 2018 が、どちらもカットオフ 0.5 ng/mL の研究で、主解析に入る 6 研究（Kato 2010、Lopez 2011、Mori 2013、Park 2016、Qureshi 2018、Rossi 2019）の中にあります。この 6 研究の要約点は感度 0.746、特異度 0.757 でした。高リスクの 2 研究を除いた 4 研究では、感度 0.709、特異度 0.781 でした（第 14 章）。部分的検証バイアスのある研究は感度が高く出やすいので、除いたときに感度が下がったことは予想と合っています。ただし、2 研究を除くと研究が 4 つに減り、区間も広がります。2 つの結果を並べて、除いても結論が変わらないかを見るために使い、どちらかを正しい結果として選ぶためには使いません。

適用可能性の懸念も同じ図や表に並べます。図 12.1 の演習では、10 研究の本文に対象の重症度や測定試薬の情報を設けていないので、適用可能性の懸念は載せていません。実際のレビューでは、適用可能性の列を加えます。例えば、ICU で術後の患者だけを対象にした研究は、バイアスのリスクが低くても、救急外来の発熱患者という問いには当てはまりにくいかもしれません。こうした研究が多いと、要約点がどの場面の精度を表しているのかがはっきりしなくなります。

\section{\hatten QUADAS-3 の変更点}
\label{sec:12-9}

2026 年 2 月に、QUADAS-2 の改訂版である QUADAS-3 が公表されました（Whiting 2026）。Cochrane DTA ハンドブック第 2 版（2023 年）と、これまでに発表されたレビューの大半は QUADAS-2 を使っているので、本書の演習は QUADAS-2 で行いました。これからレビューを始める人は、QUADAS-3 の原典を確認して、どちらを使うかをプロトコルに書いてください。

公表された論文によると、QUADAS-3 は次の 6 つの段階で評価を進めます。統合の問いを定める、その問いに対する理想的な検査精度試験を思い描く、フロー図を描く、評価する精度の推定値を特定する、バイアスのリスクと適用可能性を評価する、全体の判定を下す、の順です。バイアスのリスクの領域は、参加者、対象検査、標的疾患、解析の 4 つで、最初の 3 つでは適用可能性も評価します。

QUADAS-2 からの主な変更点として、論文は次のものを挙げています。理想的な検査精度試験という考え方を取り入れ、それと比べて実際の研究がどうずれているかを見るようにしたこと、研究単位ではなく推定値単位で評価するようにしたこと、「フローとタイミング」の領域を新しい「解析」の領域に置き換えたこと、判定の理由を書く欄を設けたこと、全体の判定を正式に導入したことです。推定値単位の評価になったことで、1 本の研究が複数のカットオフや複数の対象集団の推定値を報告している場合に、それぞれを別に評価できます。

\section{データを R に読み込む}
\label{sec:12-10}

抽出した 2×2 表は、1 研究 1 行の表にして CSV で保存します。R への読み込み方は\kkref{sec:04-4}と同じです。列は、研究名、場面、カットオフ、TP、FP、FN、TN とし、QUADAS-2 の判定も列（q\_pat、q\_idx、q\_ref、q\_flow）として加えておくと、\kkref{sec:14-5}の感度分析にそのまま使えます。判定は「低」「不明」「高」の文字で入れ、例題の 10 研究は図 12.1 の判定を入れてあります。列の名前は TP、FP、FN、TN と大文字で付けておくと、第 13 章で使う mada パッケージが自動で読み取ります。

\begin{rcode}
pct <- read.csv("data/pct_bact.csv",
                fileEncoding = "UTF-8")
names(pct)        # 列の名前を確かめる
nrow(pct)         # 研究の数（10）
\end{rcode}

\begin{rout}
 [1] "study"   "setting" "cutoff"  "TP"      "FP"      "FN"      "TN"
 [8] "q_pat"   "q_idx"   "q_ref"   "q_flow"
[1] 10
\end{rout}

読み込んだら、各研究の TP と FN の和（病気のある人の数）と、FP と TN の和（病気のない人の数）を計算し、論文の人数と一致するかを確かめます。例題では、Lopez 2011 の行で 35 + 27 = 62 と 93 + 327 = 420 になり、本文の人数と一致します。

\section{この章のまとめ}
\label{sec:12-11}

本文から 2×2 表を作るときは、人数と割合から 4 つの数を復元し、感度と特異度を計算し直して論文の値に戻ることを確かめます。複数のカットオフ、判定不能、複数のエピソードの扱いは、データを見る前にプロトコルで決めます。QUADAS-2 では 4 つの領域ごとにバイアスのリスクと適用可能性を判定し、点数にして合計しません。測定法、検体の取り扱い、採血の時点は、検査技師の目で丁寧に記録します。検査の研究に特有のバイアスは表 12.1 で整理し、初めは上段の 4 つを押さえれば足ります。判定は研究を除くためではなく、感度分析と GRADE に使います。QUADAS-2 と、2026 年公表の改訂版 QUADAS-3 のどちらを使うかは、プロトコルに書きます。次の第 13 章では、読み込んだ 10 研究に二変量モデルを当てはめます。
